\section{Joint-corner and short-shift rows}\label{sec:joint}

The two \(Q2\) populations are disjoint. This is because the chart rule
\eqref{eq:old:2_3} assigns each source to exactly one of \(u\le t\) or
\(t<u\). In a dyadic sector for the selected corners, write
\[
J\asymp J_0,\qquad K\asymp K_0,\qquad
 P_0=J_0K_0,\qquad \mu_{\rm sec}=\min(J_0,K_0).
\]
Here each dyadic sector means \(J_0\le J<2J_0\) and
\(K_0\le K<2K_0\). These half-open sectors are disjoint. The \(Q2\) predicate gives
\(\mu_{\rm sec}<Y^{1/3-j_0}\).

\begin{lemma}[Sector count]\label{lem:q2-sector}
Fix a positive selected shift \(s\) and a nonempty sector as above. Count
the canonical-source coefficient-corner triples with
\(\gcd(f,K)=\gcd(g,J)=\gcd(s,JK)=1\), with the second coefficient
forced by the conic. Suppose that
\[
fJ^2\le H,\qquad gK^2\le H.
\]
Then their number is
\[
\ll_\eta Y^\eta H\left(\frac1{P_0}+\frac{\mu_{\rm sec}^3}{P_0^2}\right)
\]
for every \(\eta>0\), with an implied constant depending only on \(\eta\).
In the \(Q2u\) application take \(H=uY\). In the \(Q2t\) application
we assert the natural cutoff \(H\ll tY^{1+\epsilon_0+l_0}\) only on
the first-match population. There the earlier \(O_R\) and \(Q1\) rows
have failed. That source-uniform bound is proved in \eqref{eq:old:5_2}.
\end{lemma}

\begin{proof}
In the \(u\) chart the conic is \(gK^2-fJ^2=u\). The imposed cutoff
gives \(fJ^2,gK^2\le H\). For the \(Q2u\) application, the source scale
identities permit the choice \(H=uY\).
The full source filter gives \(\gcd(f,K)=\gcd(g,J)=1\) and
\(\gcd(u,JK)=1\). These are the unit conditions for either orientation.
Thus the quadratic congruence has at most \(2^{\omega(K)+2}\) roots
modulo \(K^2\), or \(2^{\omega(J)+2}\) roots modulo \(J^2\).
For odd prime powers each unit-square congruence has at most two roots.
The prime \(2\) contributes at most four roots modulo a power of \(2\).
So the factor \(2^{\omega(\cdot)+2}\) is a uniform
overcount, including the even modulus.
Suppose first \(J_0\ge K_0\). For fixed \(f\) and \(K\),
\[
fJ^2\equiv-u\pmod{K^2}.
\]
This congruence has at most \(2^{\omega(K)+2}\ll_\eta Y^\eta\)
root classes modulo \(K^2\). Each class contains at most
\(J_0/K^2+1\) integers \(J\) in the sector. There are
\(O(K_0)\) choices of \(K\) and \(O(H/J_0^2)\) choices of \(f\).
The resulting count is
\[
\ll_\eta Y^\eta H
\left(\frac1{J_0K_0}+\frac{K_0}{J_0^2}\right)
=Y^\eta H\left(\frac1{P_0}+\frac{\mu_{\rm sec}^3}{P_0^2}\right).
\]
If \(K_0\ge J_0\), reduce the conic modulo \(J^2\). This gives the
same expression with \(J_0,K_0\) interchanged.

For the \(t\) chart the full filter gives
\(\gcd(g,J)=\gcd(f,K)=1\) and \(\gcd(t,JK)=1\). Use
\(gK^2-fJ^2=t\), now with \((J,K)=(B,C)\).
Combine the scale identities in \eqref{eq:old:2_2}, the source bounds
\eqref{eq:old:2_7} and the failed \(Q1\) inequality. They give
\[
fJ^2,\ gK^2\ll tY^{1+\epsilon_0+l_0}.
\AutoTag{eq:old:5_2}
\]
The source identity \(Zcd=n(n+1)\) and \(Z\ge R>Y^{5/3-\epsilon_0}\)
give
\[
cd=\frac{n(n+1)}Z<\kappa_Y Y^{1/3+\epsilon_0}.
\]
The denominator satisfies
\[
cd=ab+L\ll Y^{1/3+\epsilon_0},\quad
n-m=\frac{L\,n(n+1)}{cd(m+n+1)}
\gg Y^{1-\epsilon_0-l_0}.
\]
The second estimate uses \(n>Y/2\), \(L\ge Y^{1/3-l_0}\),
\(m+n+1\le2Y\), and the first estimate. Substitution in the last two
scale identities of \eqref{eq:old:2_2} proves \eqref{eq:old:5_2}. The same
congruence argument applies, using \(\delta C^2\equiv t\pmod{B^2}\) or
\(\gamma B^2\equiv -t\pmod{C^2}\). It gives the asserted sector bound.
\end{proof}

\begin{proposition}[Joint-corner rows]\label{prop:q2}
For every \(\varepsilon>0\),
\[
\begin{aligned}
\#Q2u&\ll_\varepsilon
Y^{2/3-2222179/15000000+\varepsilon}
+Y^{2/3-30869/1250000+\varepsilon},\\
\#Q2t&\ll_\varepsilon
Y^{2/3-972179/7500000+\varepsilon}
+Y^{2/3-S_0+\varepsilon}.
\end{aligned}\AutoTag{eq:old:5_3}
\]
The terms in each line are additive contributions to one chart
population. They are not separate first-match rows.
\end{proposition}

\begin{proof}
On \(Q2u\) or \(Q2t\), \eqref{eq:old:2_7} and the failed earlier tests give
\[
\begin{gathered}
Y^{q_{\min}}<M<\kappa_Y^2Y^{q_{\max}},\qquad
Y^{1/3-l_0}\le L<\kappa_YY^{1/3+\epsilon_0},\\
1\le s<\sqrt{\kappa_Y}\,Y^{1/6+\epsilon_0/2}.
\end{gathered}\AutoTag{eq:old:5_4}
\]
The last inequality follows from \(ut=L\) and
\(s=\min(u,t)\le\sqrt L\).

We also need a lower bound for the product of the selected corners.
Again \(Zcd=n(n+1)\) and \(Z\ge R>Y^{5/3-\epsilon_0}\), so
\(cd<\kappa_Y Y^{1/3+\epsilon_0}\). The identity
\[
\frac{n-m}{L}=\frac{n(n+1)}{cd(m+n+1)}
\]
and \(n>Y/2\) imply
\[
AD=u\frac{n-m}{L}\gg uY^{2/3-\epsilon_0}.
\]
Reducing the exponent by \(\tfrac32\epsilon_0\) absorbs the fixed
constant for all sufficiently large \(Y\). A sector has
\(AD<4P_0\). So its nonemptiness implies
\[
P_0>\frac{uY^{2/3-5\epsilon_0/2}}{32}
\quad\text{in the \(u\) chart.}\AutoTag{eq:old:5_5}
\]
In the \(t\) chart, \(BC=t(m+n+1)/L\), so \eqref{eq:old:5_4} gives
\[
BC\gg tY^{2/3-\epsilon_0},\qquad
P_0>\frac{tY^{2/3-5\epsilon_0/2}}{32}.
\AutoTag{eq:old:5_6}
\]
We reduce the exponent by \(3\epsilon_0/2\) from the scale bound
\(BC\gg tY^{2/3-\epsilon_0}\). This absorbs fixed dyadic constants uniformly.
This gives the exponent used in \eqref{eq:old:5_8}.

Apply Lemma~\ref{lem:q2-sector} in each occupied sector. In the \(u\)
chart, its two terms are bounded by \eqref{eq:old:5_5} and
\(\mu_{\rm sec}<Y^{1/3-j_0}\) as
\[
Y^\eta\left(
Y^{1/3+5\epsilon_0/2}
+\frac{Y^{2/3-3j_0+5\epsilon_0}}u
\right).
\AutoTag{eq:old:5_7}
\]
For the \(t\) chart, \eqref{eq:old:5_2}, \eqref{eq:old:5_6}, and the same corner bound give
\[
Y^\eta\left(
Y^{1/3+7\epsilon_0/2+l_0}
+\frac{Y^{2/3-3j_0+6\epsilon_0+l_0}}t
\right).
\AutoTag{eq:old:5_8}
\]
The source has \(s\le\sqrt{\kappa_Y}Y^{1/6+\epsilon_0/2}\).
Sum the first terms of \eqref{eq:old:5_7}--\eqref{eq:old:5_8} over this many
shifts. This adds \(1/6+\epsilon_0/2\) to their exponents. The second terms are harmonic
in the selected shift and add only \(O(\log Y)\). There are
\(O(\log^2Y)\) corner sectors.

Fix a coefficient \(f\), a sector corner pair and a shift. Then the
conic forces the other coefficient \(g\). The number of ordered residual
splits of \(f,g\) is at most \(\tauDiv(f)\tauDiv(g)\ll_\eta Y^\eta\).
The inverse maps \eqref{eq:old:2_2} then determine the remaining corners. The full
canonical filter is retained. Choose each internal exponent first. Then
the divisor factors, root factors, dyadic logarithms and the harmonic sum
use a total arbitrarily small loss.

Substitute \(\epsilon_0=k_0=S_0\), \(l_0=2S_0\) and \(j_0=3S_0\).
This gives the four exponents
\[
\begin{array}{ll}
1/2+3S_0=2/3-2222179/15000000,&
2/3-9S_0+5S_0=2/3-30869/1250000,\\
1/2+4S_0+2S_0=2/3-972179/7500000,&
2/3-9S_0+6S_0+2S_0=2/3-S_0.
\end{array}
\AutoTag{eq:old:5_9}
\]
This proves \eqref{eq:old:5_3}.
\end{proof}
